\section{Methodology}
\label{sec:method}

\subsection{Ranging and range calibration}
Each range is obtained by a double-sided two-way exchange (DS-TWR),
whose symmetric four-timestamp formulation cancels the clock-frequency
offset between the two uncoordinated crystals to first
order~\cite{dstwr_neirynck} (Fig.~\ref{fig:dstwr}). One DW1000 timer
tick corresponds to \SI{4.69}{\milli\meter} of range and sets the
resolution of every calibration constant below. The two systematic
biases in~\eqref{eq:measmodel} are removed in sequence. The
antenna-delay term $\delta_k$, which reflects the propagation time
internal to each board, is calibrated per anchor against a reference
position and stored on the device, after which all anchors agree with
the reference distance to within $\pm\SI{16}{\milli\meter}$. The
residual power-dependent bias $\gamma_k$ is modelled per anchor as an
affine function of the first-path power $P_{\mathrm{fp}}$, which is more
stable than total power,
\begin{equation}
  \hat{d}_k = \rho_k - \big(a_k\,P_{\mathrm{fp},k} + b_k\big),
  \qquad a_k\in[7.7,\,9.5]\ \si{\milli\meter\per\dB},
  \label{eq:power}
\end{equation}
with $P_{\mathrm{fp}}$ clamped to the fitted interval so that the model
is never extrapolated. The coefficients are estimated once by
leave-one-position-out regression over a static campaign.

\begin{figure}[t]
  \centering
  \begin{tikzpicture}[font=\small, >={Stealth[length=2.2mm]}, yscale=0.9]
    \node[font=\bfseries, text=cBlue]  at (0,0.4)  {Tag};
    \node[font=\bfseries, text=cAmber] at (4.6,0.4){Anchor};
    \draw[cBlue!55, line width=1pt]  (0,0.1)   -- (0,-5.4);
    \draw[cAmber!70, line width=1pt] (4.6,0.1) -- (4.6,-5.4);
    \draw[->, line width=1pt, draw=cBlue]  (0,-0.3)  -- node[above,sloped,font=\scriptsize]{Poll}     (4.6,-0.9);
    \draw[->, line width=1pt, draw=cAmber] (4.6,-2.0) -- node[above,sloped,font=\scriptsize]{Response} (0,-2.6);
    \draw[->, line width=1pt, draw=cBlue]  (0,-3.7)  -- node[above,sloped,font=\scriptsize]{Final}    (4.6,-4.3);
    \foreach \y/\t in {-0.3/t_1, -2.6/t_2, -3.7/t_3} \node[left=1pt,font=\scriptsize] at (0,\y) {$\t$};
    \foreach \y/\t in {-0.9/t_4, -2.0/t_5, -4.3/t_6} \node[right=1pt,font=\scriptsize] at (4.6,\y) {$\t$};
    \draw[decorate,decoration={brace,amplitude=4pt},cGray]
      (-0.75,-0.3) -- node[left=4pt,font=\scriptsize]{$T_{\mathrm{round}1}$} (-0.75,-2.6);
    \draw[decorate,decoration={brace,amplitude=4pt},cGray]
      (-0.75,-2.6) -- node[left=4pt,font=\scriptsize]{$T_{\mathrm{reply}2}$} (-0.75,-3.7);
    \draw[decorate,decoration={brace,amplitude=4pt,mirror},cGray]
      (5.35,-0.9) -- node[right=4pt,font=\scriptsize]{$T_{\mathrm{reply}1}$} (5.35,-2.0);
    \draw[decorate,decoration={brace,amplitude=4pt,mirror},cGray]
      (5.35,-2.0) -- node[right=4pt,font=\scriptsize]{$T_{\mathrm{round}2}$} (5.35,-4.3);
  \end{tikzpicture}
  \caption{Double-sided two-way ranging. Four host-side timestamps over
  two round trips recover the time-of-flight without a shared clock; the
  symmetric form is first-order insensitive to crystal frequency offset.}
  \label{fig:dstwr}
\end{figure}

\subsection{Robust NLOS-aware estimation}
Under obstruction the first signal component is attenuated relative to
the total received energy, so the gap between total and first-path power
is a direct indicator of the NLOS excess $\Delta_{ik}$
(Fig.~\ref{fig:nlos}). Denoting this gap $g_k = P_{\mathrm{rx},k} -
P_{\mathrm{fp},k}$ in decibels, each range is assigned a continuous
weight
\begin{equation}
  w_k = \max\!\Big(w_{\mathrm{f}},\ 10^{-\max(0,\,g_k-\tau)/10}\Big),
  \label{eq:nlos}
\end{equation}
with LOS threshold $\tau=\SI{3}{\dB}$ and floor $w_{\mathrm{f}}=0.05$:
below the threshold a range keeps unit weight, above it the weight decays
with the excess gap, and the floor prevents any single anchor from being
discarded outright. The instantaneous fix follows from robust weighted
multilateration,
\begin{equation}
  \pvec^\star = \arg\min_{\pvec}\;
  \sum_{k=1}^{N} w_k\,\rho_\delta\!\big(\lVert \pvec-\avec_k\rVert - \hat{d}_k\big),
  \label{eq:ls}
\end{equation}
where $\rho_\delta$ is the Huber loss and a $\chi^2$ consistency gate
rejects gross outliers, so that an inconsistent range bends rather than
dominates the solution.

The trajectory itself is estimated by a moving-horizon estimator (MHE)
that, at each step, jointly optimizes the last $N_{\mathrm{h}}$ states
($N_{\mathrm{h}}{=}10$, $\approx\!\SI{2}{\second}$) so that new evidence
can revise the recent estimate:
\begin{equation}
\begin{aligned}
  \min_{\{\xvec_k\}}\ &
  \sum_{k}\sum_{j} w_{kj}\,\rho_\delta\!\big(\lVert\pvec_k-\avec_j\rVert - \hat{d}_{kj}\big)\\[-1pt]
  &{}+ \sum_{k}\big\lVert \xvec_{k+1}-f(\xvec_k)\big\rVert^2_{Q^{-1}}
  + \big\lVert \xvec_1-\bar{\xvec}\big\rVert^2_{P^{-1}},
\end{aligned}
\label{eq:mhe}
\end{equation}
where the first term fits every windowed range under the same weighted
robust loss, the second penalizes departures from the kinematic motion
model $f(\cdot)$, and the third is an arrival cost coupling the window to
the previous estimate. The same weights $w_{kj}$ scale the measurement
variances, so the NLOS score governs both the fix and the filter. The
IMU contributes only differential quantities---yaw rate and
gravity-referenced roll and pitch---never an absolute heading; because
the Ackermann platform cannot slip laterally, $f(\cdot)$ constrains the
body velocity to the heading direction. For two tags, the rigid
constraint~\eqref{eq:rigid} is imposed by construction, so a single body
state generates both tag predictions, the baseline is exact in every
solution, and the heading is observed directly by the radios with a
per-reading resolution $\sigma_\psi \approx \sqrt{2}\,\sigma_d/L$
($\approx\!4.6^\circ$ for $\sigma_d\!=\!30$\,mm, $L\!=\!0.527$\,m),
further reduced by filtering.

\begin{figure}[t]
  \centering
  \begin{tikzpicture}[font=\small, >={Stealth[length=2.2mm]}]
    \draw[cGray!50, line width=1pt] (-0.3,-1.5) -- (7.1,-1.5);
    \node[cGray, font=\scriptsize] at (6.7,-1.7) {floor};
    \fill[cGray!25, draw=cGray!70] (3.3,-1.5) rectangle (3.9,1.15);
    \node[cGray!80!black, font=\scriptsize\bfseries, align=center] at (3.6,1.5) {obstacle};
    \node[circle, draw=cBlue, fill=cBlueL, line width=1pt, inner sep=2.5pt]
      (tag) at (0,0) {\footnotesize\bfseries Tag};
    \node[rectangle, rounded corners=2pt, draw=cAmber, fill=cAmberL,
      line width=1pt, inner sep=3.5pt] (anc) at (6.6,0) {\footnotesize\bfseries Anchor};
    \draw[dashed, cGray, line width=1pt] (tag) -- (anc);
    \node[cRed, font=\bfseries] at (3.6,0) {$\times$};
    \node[cRed, font=\scriptsize\bfseries, align=center] at (2.05,0.32) {direct path\\ blocked};
    \draw[cGreen, line width=1.4pt, rounded corners=6pt]
      (tag.south) -- (1.7,-1.15) -- (3.6,-1.15) -- (5.2,-1.15) -- (anc.south);
    \node[cGreen!60!black, font=\scriptsize] at (3.6,-0.95) {reflected / diffracted};
    \node[align=center, font=\small] at (3.3,0.78) {$\rho_k = d + \Delta_k,\ \ \Delta_k>0$};
  \end{tikzpicture}
  \caption{NLOS mechanism. When the direct path is obstructed the signal
  reaches the anchor along a longer reflected or diffracted route, so the
  measured range is biased positively by $\Delta_k$. The bias is
  estimated from the total-to-first-path power gap and used to weight the
  range through~\eqref{eq:nlos}.}
  \label{fig:nlos}
\end{figure}

\subsection{Automatic anchor localization}
When the anchor coordinates are not measured externally, they are
recovered from the ten pairwise inter-anchor ranges of the five-anchor
set. A set of mutual distances defines a rigid configuration up to an
isometry, and classical multidimensional scaling~\cite{torgerson1952mds}
returns the point set whose pairwise distances best match the measured
ranges. The triggering tag only relays messages, so its own antenna
delay does not enter the survey. The recovered constellation is
determined up to a rotation and reflection, resolved against a single
known landmark or, for evaluation, aligned to a reference layout without
scaling.
